\section{Related Work}
\label{sec:related}



% \subsection{Spam and Phishing Detection}

Spam and phishing email detection has evolved through multiple paradigm shifts driven by adversarial co-evolution between attackers and defenders. Early systems relied on rule-based filtering and manually curated blacklists~\cite{al2023multi,magdy2022efficient}, which proved brittle against adaptive attackers. The field subsequently transitioned to machine learning approaches, with statistical classifiers including Naive Bayes, Support Vector Machines, and ensemble methods~\cite{kontsewaya2021evaluating,wijaya2023spam,budiman2024email,bazzaz2021hybrid,wang2025hybrid} demonstrating superior generalization through learned feature representations. Recent advances in deep learning have further improved detection capabilities, with transformer-based models leveraging contextualized embeddings~\cite{altwaijry2024advancing,thakur2023systematic}.

However, two persistent challenges limit practical deployment effectiveness. First, severe class imbalance characterizes operational environments where malicious emails constitute only 1-10\% of total traffic~\cite{bazzaz2021hybrid,johari2025key}, causing classifiers to develop strong bias toward majority classes and sacrificing minority class sensitivity~\cite{wang2025hybrid}. Second, privacy regulations and legal constraints restrict retention and sharing of email data, limiting access to large-scale training datasets~\cite{kim2023pcsf}. These dual challenges motivate exploration of synthetic data generation as a potential solution.

% \subsection{Synthetic Data Generation for Cybersecurity}

Synthetic data generation addresses training data scarcity by creating artificial examples that augment limited real datasets. Traditional approaches in natural language processing have explored augmentation techniques including synonym replacement, back-translation, and paraphrasing~\cite{shorten2021text}, though these methods lack semantic understanding and cannot generate content with controlled characteristics while maintaining realistic coherence. Deep generative models such as Generative Adversarial Networks have shown effectiveness for structured cybersecurity data including network traffic and intrusion detection~\cite{bourou2021review,ammara2025syntheticnetworktrafficdata,10765501,agrawal2024review}, but face critical limitations for text generation due to difficulties maintaining semantic coherence and linguistic realism required for spam generation.

Large Language Models(LLMs) have fundamentally transformed synthetic text generation for spam and phishing attacks through pre-training on massive corpora spanning diverse linguistic patterns. Unlike feature-space interpolation or adversarial training, LLMs generate text through autoregressive prediction conditioned on context, enabling semantically coherent content that maintains consistency across multiple attributes~\cite{razmyslovich2025eltex,almorjan2025large}. This makes LLMs the preferred choice for attackers requiring rapid, large-scale, cost-effective spam generation through simple API calls without model training. Recent applications to cybersecurity have explored domain-driven synthetic text generation for classification tasks~\cite{razmyslovich2025eltex,zhang2025llms}, structured data creation incorporating domain knowledge~\cite{almorjan2025large}, and general cyber defense applications~\cite{ammara2024synthetic}. Systematic reviews analyzing hundreds of LLM applications~\cite{zhang2025llms,ding2023tmg} document broad adoption across security domains, while emerging research explores prompt engineering strategies for controlled generation~\cite{chen2025unleashing,de2023unsupervised}.

However, existing work exhibits critical limitations. Prior evaluations focus predominantly on single-sided assessments either defensive utility for classifier training or offensive threats for spam generation without systematically characterizing the dual-use nature of LLMs where the same generation capabilities empower both attackers and defenders. Recent work has demonstrated LLMs' potential for generating convincing phishing emails and automating social engineering attacks~\cite{bethany2024large,afane2024next,weinz2025impact,ai2024defending,thapa2025large}, yet comprehensive evaluation frameworks simultaneously assessing both offensive threats and defensive countermeasures under controlled experimental conditions remain absent. Furthermore, cross-model generalization dynamics that whether training on synthetic data from one LLM improves detection of spam from different LLM sources remain unexplored despite operational importance when attackers and defenders employ different models. Rigorous experimental methodology employing replicated trials and systematic exploration of generation parameters is largely absent from prior work.
